% !TEX root = paper.tex
\doclearpage
\section{Technical Overview}
\label{sec:our_techniques}




\subsection{String Operations as Relational Operations}
\label{sec:string-ops-relational}



\parhead{Strings as relations} The triple $(\mathsf{orig\text{-}ind}, \mathsf{int\text{-}ind}, \mathsf{char})$ is an ordinary relational table with one row per character, which we call the \emph{auxiliary character table} and write $\strcol.\mathsf{aux\_chars}$; the hash column with the string index forms the \emph{main table} $\strcol.\mathsf{main}$, with one row per string, and the string-level auxiliary columns, such as the length $\mathsf{l}$, form the \emph{auxiliary string table} $\strcol.\mathsf{aux\_strs}$ over the same rows. In the auxiliary character table, $\mathsf{orig\text{-}ind}$ is a foreign key into the main table and $(\mathsf{orig\text{-}ind}, \mathsf{int\text{-}ind})$ is a primary key. Once a string is a set of $(\text{string}, \text{position}, \text{character})$ rows, the operations that are opaque on the hash column become plain SQL. 

\parhead{Pattern matching is a filter and a join} Take the pattern \texttt{a\%b}, a one-character prefix and a one-character suffix with anything in between. A string matches iff its character rows contain an \texttt{a} with $\mathsf{int\text{-}ind}=0$ and a \texttt{b} at the last position; the suffix has to know where the string ends, so we record each string's length $\mathsf{l}$ in the auxiliary string table. The two characters live in two different rows of the auxiliary character table, so the query joins two copies of it on the string they belong to, pins the \texttt{a} to the first position and the \texttt{b} to the last one, and asks that the \texttt{b} come strictly after the \texttt{a}:
\begin{equation}\label{eq:like-as-join}
\begin{array}{@{}l@{}}
\sqlline{\sqlkw{SELECT} m.ind}\\
\sqlline{\ \ \sqlkw{FROM} s.aux\_strs m, s.aux\_chars x, s.aux\_chars y}\\
\sqlline{\ \sqlkw{WHERE} x.orig\_ind = m.ind \sqlkw{AND} y.orig\_ind = m.ind}\\
\sqlline{\ \ \ \sqlkw{AND} x.int\_ind = 0 \sqlkw{AND} x.char = 'a'}\\
\sqlline{\ \ \ \sqlkw{AND} y.int\_ind = m.l - 1 \sqlkw{AND} y.char = 'b'}\\
\sqlline{\ \ \ \sqlkw{AND} y.int\_ind > x.int\_ind;}
\end{array}
\end{equation}
The general case has the same three ingredients. Each literal character is one more copy of $\strcol.\mathsf{aux\_chars}$, tied to its string by $\mathsf{orig\text{-}ind}$ and pinned to a position by $\mathsf{int\text{-}ind}$, so \texttt{ab\%} adds a copy for the \texttt{b} at position $1$ and a \_ simply skips a position; the length lets a suffix find the end of the string; and each \% becomes an inequality between the positions of consecutive factors, as on the last line. The result is always an equi-join of $\strcol.\mathsf{aux\_chars}$ with one copy per literal character, a filter on the characters, and a projection onto the string index, i.e.\ a \texttt{GROUP BY orig\_ind}, to obtain one bit per string. Matching a pattern by joining one copy of a table per symbol of the pattern is the same construction that query engines use, over rows rather than characters, to pre-filter row-pattern queries~\cite{ZhuHC23}.

\parhead{Ordering is a join followed by an aggregate} Two strings compare lexicographically at the first position where they differ. Take \texttt{apple} and \texttt{apricot}, held in rows $i$ and $j$ of the column. Joining their character rows position by position and keeping the pairs whose characters differ leaves positions $2$ (\texttt{p} vs.\ \texttt{r}), $3$ and $4$; the first of them decides:
\begin{equation}\label{eq:order-as-join}
\begin{array}{@{}l@{}}
\sqlline{\sqlkw{SELECT} \sqlkw{MIN}(x.int\_ind)}\\
\sqlline{\ \ \sqlkw{FROM} s.aux\_chars x, s.aux\_chars y}\\
\sqlline{\ \sqlkw{WHERE} x.orig\_ind = i \sqlkw{AND} y.orig\_ind = j}\\
\sqlline{\ \ \ \sqlkw{AND} x.int\_ind = y.int\_ind}\\
\sqlline{\ \ \ \sqlkw{AND} x.char <> y.char;}
\end{array}
\end{equation}
The query returns $2$, and since \texttt{p} $<$ \texttt{r} at that position, \texttt{apple} precedes \texttt{apricot}. If the query returns nothing, one string is a prefix of the other, and the length $\mathsf{l}$ in the auxiliary string table settles the order. A join, a filter and an aggregate thus compare any two strings, and sorting a column compares every pair of adjacent rows this way. \Cref{sec:rank} goes one step further and precomputes this order once, as a \emph{rank} column, so that sorting strings becomes sorting integers.

\parhead{Our Approach} The relational view shows that once the characters are exposed, every white-box operation we need is \emph{provable}, since each is a composition of filters, joins and aggregates over the three tables, and $\truthtable$ proves all of these. Our actual protocols, however, do not hand these queries to $\truthtable$ as a black box. Instead, we unroll the PIOPs that $\truthtable$ would run for them and design special-purpose PIOPs of roughly the same shape, a lookup where the query joins, a zerocheck where it filters, a keyed sumcheck where it aggregates, that additionally exploit what a generic query engine cannot know about the tables.

\subsection{Acceleration via Preprocessing}
\label{sec:preprocessing-overview}

Although, the above approach gives us complete and sound protocols, they are rather complex. For example, for a filter as simple as \texttt{s LIKE '\%na\%'}, the protocol of \cref{sec:piop-for-pattern-matching} issues about $24$ zerocheck claims, $13$ lookup claims, $4$ keyed-sumcheck claims and $4$ nonzero-check claims, most of them over the bloated character-level hypercube, i.e.\ with one entry per character rather than per string, and every further factor in the pattern adds about $15$ zerochecks and $10$ lookups. This is because reasoning about characters are inherently complex and depends on the length of the strings which makes our protocols hard to scale.

To mitigate this problem, we introduce two categories of preprocessing techniques and leverage the simple fact that since data owner is trusted, it can provide more useful information at the commitment time which are \emph{query-independent}.

\parhead{Preprocessed Ranking} We first start with the simple idea of preprocessed rankings which makes proving typical string sorting queries as cheap as it gets. The observation is that the strings a query sorts are almost never new. Relational operators such as filters and joins only copy strings around; a new string arises only from a string function such as \texttt{SUBSTRING} or \texttt{CONCAT}, and such sort keys are rare: of the ten TPC-H queries~\cite{PoessF00} that order by a string column, nine sort a base column or a renamed copy of one, and only one sorts a computed string, a two-character prefix of a phone number. So each sorted string is, with rare exceptions, a copy of a string in a committed base column, and the \emph{relative order} of those strings is fixed by the base column alone. The data owner can therefore compute this order once, at ingestion, and store it next to each string as an integer. When the sort finally happens, the prover sorts the rank column instead of the strings, which is exactly the integer sort that $\truthtable$ already proves; the characters are never touched, so the cost of a string sort no longer depends on the lengths of the strings.
As an example, take a table $T$ whose string column $\strcol$ is the one of \cref{fig:string-arith-example}, and the query on the left, which we rewrite into the query on the right:
\begin{equation}\label{eq:order-by-rank}
\begin{array}{@{}l@{}}
\sqlline{\sqlkw{SELECT} s}\\
\sqlline{\ \ \sqlkw{FROM} T}\\
\sqlline{\ \sqlkw{ORDER BY} s;}
\end{array}
\quad\Longrightarrow\quad
\begin{array}{@{}l@{}}
\sqlline{\sqlkw{SELECT} s}\\
\sqlline{\ \ \sqlkw{FROM} T}\\
\sqlline{\ \sqlkw{ORDER BY} s.rk;}
\end{array}
\end{equation}
Here \texttt{s.rk} is the rank column of $\strcol$, which in \cref{fig:string-arith-example} holds $0$ for \texttt{apple}, $1$ for \texttt{banana}, $2$ for \texttt{cherry} and $3$ for \texttt{orange}. Running the right query sorts the active rows by these integers, $0, 1, 3$, and returns \texttt{apple}, \texttt{banana}, \texttt{orange}: exactly the result of the left query, obtained without comparing a single character, and much simpler and faster.

\parhead{Preprocessed Fingerprints} Pattern matching has no analogue of the rank, since the pattern is only known at query time, but it admits a cheap \emph{pre-filter}. 
The idea is to first run a cheap filter that has false positives but no false negatives; i.e. it discards most of the rows that are supposed to be dropped and never discards a match, so the expensive protocol then runs on a much smaller table. This is exactly the trade-off of a Bloom filter~\cite{Bloom70}, which answers membership queries with false positives but no false negatives, and our filter can be seen as one.

The filter, borrowed from the database literature~\cite{string_fingerprint}, works as follows. Consider running the pattern \texttt{\%na\%} on the string column of \cref{fig:string-arith-example}; clearly, \texttt{banana} is the only match. Now consider simple \emph{features} of a string, such as containing the character \texttt{n} or containing the character \texttt{a}. The pattern \texttt{\%na\%} has both features, so any string that could possibly match must have them too, and a string that lacks either can be dropped without looking any further: \texttt{apple} contains no \texttt{n} and is gone at once. \texttt{orange}, on the other hand, has both characters and survives as a false positive. Still, instead of running the expensive protocol on all four rows, we now run it on the two that survived. Our pre-filtering rests on the observation that most strings can be dropped this way, simply by checking that they lack a feature the pattern has, which shrinks the search space for the actual matches. 

In general, we fix a list of $m$ features, such as containing a given character or a given short run of characters (an $n$-gram). The \emph{fingerprint} of a string is a bit vector with one bit per feature, set iff the string has that feature; the fingerprint of a pattern is computed in the same way from its literal characters, separately from the data. Since the fingerprints of the strings depend only on the data, the data owner preprocesses them at commitment time: it computes the fingerprint of every string once and commits the fingerprints alongside the column, so that at query time only the pattern's fingerprint remains to be computed. For all this to pay off, the pre-filtering must come with a protocol that is substantially cheaper than the actual pattern-matching protocol; otherwise it is pointless. We devise such a protocol in \cref{sec:fp-piop}.

